% !TEX root =../ECE343_Slides_Rogers.tex

%%%%%%%%%%%%%%%%%%%% LO 3.3%%%%%%%%%%%%%%%%%
%%%%%%%%%%%%%%%%%%%%%%%%%%%%%%%%%%%%%%%%%%%%%%%%%%%%%%
%%%%%%%%%%%%%%%%%%%%%%%%%%%%%%%%%%%%%%%%%%%%%%%%%%%%%%
%\begin{comment}
\section{Schedule}
\renewcommand{\arraystretch}{1}
\begin{frame}{}
	\begin{textblock}{14}(0.5,1.75)
	\def\myfontsize{footnotesize}
	
\setcounter{mylnum}{1}
\begin{center}
\begin{\myfontsize}
\begin{tabular}{| c | l | c | l |}
\hline
\textbf{Lesson} 		  	& \textbf{Topics}             & \pbox{5cm}{\textbf{Reading}}	& \textbf{Notes}		\\ \hline\hline
\ifnum\lessonNum=\value{mylnum} \rowcolor{green!20!white}\fi
	\themylnum  
	& Course Overview  \& Introduction		
	& Ch.1 
	&  LO 1.1-1.2 
	\stepcounter{mylnum}\\ \hline \hline
\rowcolor{gray!20!white}\multicolumn{4}{| c |}{\textbf{Chapter 2: Electric and Magnetic Fields}}	\\ \hline\hline
\ifnum\lessonNum=\value{mylnum} \rowcolor{green!20!white}\fi
	\themylnum 	
	& Fields, Parameters, Field Quantities	
	& Ch. 2	
	& LO 2.1-2.2			
\stepcounter{mylnum}\\ \hline\hline
\rowcolor{gray!20!white}\multicolumn{4}{| c |}{\textbf{Chapter 3: Transmission Lines}}	 \\ \hline\hline	
\ifnum\lessonNum=\value{mylnum} \rowcolor{green!20!white}\fi
	\themylnum 	
	& Lumped Element Model				
	& 3.1-3.4 
	& LO 3.1		
\stepcounter{mylnum}\\ \hline
\ifnum\lessonNum=\value{mylnum} \rowcolor{green!20!white}\fi
	\themylnum 	
	& Telegrapher's Equations
	& 3.5-3.6
	&  LO 3.1	
\stepcounter{mylnum}\\ \hline
\ifnum\lessonNum=\value{mylnum} \rowcolor{green!20!white}\fi
	\themylnum 	
	&  Characteristic Impedance \& Wave Propagation	
	& 3.7-3.9		
	& LO 3.2						
\stepcounter{mylnum}\\ \hline
\ifnum\lessonNum=\value{mylnum} \rowcolor{green!20!white}\fi
	\themylnum 	
	& Coaxial,  Reflection Coefficient, VSWR	
	& 3.10,3.12-3.14	
	& LO 3.2-3.3		
\stepcounter{mylnum}\\ \hline 
\ifnum\lessonNum=\value{mylnum} \rowcolor{green!20!white}\fi
	\themylnum 	
	&  Terminated Transmission Lines	
	&  3.15,3.17-3.19
	& LO 3.4		
\stepcounter{mylnum}\\ \hline 
\ifnum\lessonNum=\value{mylnum} \rowcolor{green!20!white}\fi
	\themylnum	
	& Impedance Matching	
	& 3.21-3.23
	& LO 3.4 			
\stepcounter{mylnum}\\ \hline 
\ifnum\lessonNum=\value{mylnum} \rowcolor{green!20!white}\fi
	\themylnum  
	& Impedance Matching	
	& 3.21-3.23
	& LO 3.4		
\stepcounter{mylnum}\\ \hline 
\ifnum\lessonNum=\value{mylnum} \rowcolor{green!20!white}\fi
	\themylnum	
	& Smith Charts
	& Supplemental
	& LO 3.5
\stepcounter{mylnum}\\ \hline 
\ifnum\lessonNum=\value{mylnum} \rowcolor{green!20!white}\fi
	\themylnum	
	& Smith Charts
	& Supplemental 
	& LO 3.5					
\stepcounter{mylnum}\\ \hline 
\ifnum\lessonNum=\value{mylnum} \rowcolor{green!20!white}\fi
	\themylnum	
	& Smith Charts
	& Supplemental 
	& LO 3.5	
\stepcounter{mylnum}\\ \hline 
\ifnum\lessonNum=\value{mylnum} \rowcolor{green!20!white}\fi
	\themylnum	
	& Smith Charts
	& Supplemental 
	& LO 3.5	
\stepcounter{mylnum}\\ \hline 

\ifnum\lessonNum=\value{mylnum} \rowcolor{green!20!white}\fi
	\themylnum	
	& Transients		
	& Supplemental 
	& LO 3.6	 
\stepcounter{mylnum}\\ \hline \hline
\ifnum\lessonNum=\value{mylnum} \rowcolor{green!20!white}\fi
	\themylnum 	
	& Flex Day
	& 
	& 	
\stepcounter{mylnum}\\ \hline 
\ifnum\lessonNum=\value{mylnum} \rowcolor{green!20!white}\fi
	\themylnum 	
	& Reflection/Transmission Coefficient Lab
	& Lab Packet
	& LO 3.7		
\stepcounter{mylnum}\\ \hline 
\ifnum\lessonNum=\value{mylnum} \rowcolor{green!20!white}\fi
	\themylnum 	
	& Reflection/Transmission Coefficient Lab
	& Lab Packet
	& LO 3.7		
\stepcounter{mylnum}\\ \hline 
\rowcolor{gray!20!white}\multicolumn{4}{| c |}{\textbf{Chapter 4: Vector Analysis}} \\ \hline\hline
\ifnum\lessonNum=\value{mylnum} \rowcolor{green!20!white}\fi
	\themylnum	
	& Vector \& Coordinate System Review
	& Ch. 4
	&LO 4.1-4.4 \textbf{GR 1 due}	
\stepcounter{mylnum}\\ \hline 
\rowcolor{gray!20!white}\multicolumn{4}{| c |}{\textbf{Chapter 5: Electrostatics}}\\ \hline\hline
\ifnum\lessonNum=\value{mylnum} \rowcolor{green!20!white}\fi
	\themylnum	
	& Coulomb's Law
	& 5.1-5.3
	& LO 5.1
\stepcounter{mylnum}\\ \hline 
\ifnum\lessonNum=\value{mylnum} \rowcolor{green!20!white}\fi
	\themylnum	
	& Surface and Volume Charge			
	& 5.4
	& LO 5.1
\stepcounter{mylnum}\\ \hline 
\ifnum\lessonNum=\value{mylnum} \rowcolor{green!20!white}\fi
	\themylnum	
	& Gauss' Law (Integral Form)
	& 5.5-5.6
	& LO 5.2		 					
\stepcounter{mylnum}\\ \hline 
\ifnum\lessonNum=\value{mylnum} \rowcolor{green!20!white}\fi
	\themylnum	
	& Gauss' Law (Differential Form)		
	& 5.7-5.8
	& LO 5.2	 	
\stepcounter{mylnum}\\ \hline 
\ifnum\lessonNum=\value{mylnum} \rowcolor{green!20!white}\fi
	\themylnum	
	& Electric Potential	
	& 5.9-5.10, 5.15
	&  LO 5.3	
\stepcounter{mylnum}\\ \hline 
% <---- line 129 should be 20 lessons, based on 6 lines per lesson, starting at line 9 
\end{tabular}			
\end{\myfontsize}
\end{center}

	\end{textblock}
\end{frame}

%%%%%%%%%%%%%%%%%%%%%%%%%%%%%%%%%%%%%%%%%%%%%%%%%%%%%
%%%%%%%%%%%%%%%%%%%%%%%%%%%%%%%%%%%%%%%%%%%%%%%%%%%%%

\def\obja{I can describe the origin of the transmission line equation, understand the method of developing a solution, and interpret the solution.}
\def\objb{I can explain the physical meaning of and calculate the following quantities: propagation constant, phase constant, attenuation constant, characteristic impedance, VSWR, reflection coefficient, and transmission coefficient.}
\def\objc{I can use the transmission line equation to calculate the voltage/current at an arbitrary point on a specified transmission line.}
\def\objd{ I can use the transmission line equation to match a load to a source using a length of transmission line.}
\def\obje{ I can identify the various quantities represented on the Smith Chart.}
\def\objf{ I can use a Smith Chart to design a matching network for a given source and load, using series components, parallel components, shorted stubs, and open stubs.}
\def\objg{ I can plot the transient response of a load connected to a source through a transmission line. }
\def\objh{ I can measure the reflection coefficient across a range of frequencies and interpret the measurement using a network analyzer.}
\begin{frame}{Ch 3 Learning Objectives}
	\begin{textblock}{14}(0.5,3)
	\begin{scriptsize}
	\begin{enumerate}[{3}.1]
		\ifnum\objNum=1
			\item \alert{ \obja }
		\else 
			\item \obja
		\fi 
		\ifnum\objNum=23
			\item \alert{ \objb }
			\item \alert{ \objc }
		\else 
			\item \objb
			\item \objc
		\fi 
%		\ifnum\objNum=3
%			\item \alert{ \objc }
%		\else 
%			\item \objc
%		\fi 
		\ifnum\objNum=4
			\item \alert{ \objd }
		\else 
			\item \objd
		\fi 
		\ifnum\objNum=56
			\item \alert{ \obje }
		\else 
			\item \obje
		\fi 
		\ifnum\objNum=56
			\item \alert{ \objf }
		\else 
			\item \objf
		\fi 
		\ifnum\objNum=7
			\item \alert{ \objg }
		\else 
			\item \objg
		\fi 
		\ifnum\objNum=8
			\item \alert{ \objh }
		\else 
			\item \objh
		\fi 

	\end{enumerate}
	\end{scriptsize}
	\end{textblock}
\end{frame}

%%%%%%%%%%%%%%%%%%%%%%%%%%%%%%%%%%%%%%%%%%%%%%%%%%%%%%%%%%%%%%%%%
%%%%%%%%%%%%%%%%%%%%%%%%%%%%%%%%%%%%%%%%%%%%%%%%%%%%%%%%%%%%%%%%%
% LO 3.5
\section{Smith Charts}
\begin{frame}{Smith Charts: History}
	\begin{textblock}{14.5}(0.5,3)

	\begin{itemize}
		\setlength\itemsep{18pt}
		\item Developed in 1939 by Phillip H. Smith
		\item Used to solve a variety of tx line and waveguide problems
		\item  Measure/evaluate complex reflection/transmission on a tx line
		\item Measure/evaluate VSWR
		\item Measure/evalute impedance along a tx line
		\item Evaluate IO transfer functions
	\end{itemize}
	
	\end{textblock}
	
\end{frame}

%%%%%%%%%%%%%%%%%%%%%%%%%%%%%%%%%%%%%%%%%%%
%%%%%%%%%%%%%%%%%%%%%%%%%%%%%%%%%%%%%%%%%%%
\begin{frame}{Smith Chart}
\begin{textblock}{14.5}(0.5,2)
\begin{center}
\begin{tikzpicture}[scale=0.35]

\input{\mygraphicspath/ECE343/smithChartFoundation}

   \begin{smithchart}[
                      show origin,
                      width=20cm,
                     ]
%   \addplot[mark=none,line width=2]
%       coordinates{
%           (1, 0) (1, 0.1) (1,0.2) (1,0.3) (1,0.4) (1,0.5) (1,0.5)
%       };
%   \addplot[mark=none,line width=1, red]
%       coordinates{
%           %(0,-0.5)(1, 0) (0, 0.5) % pacman angle
%           (1,0) (0.5,0.5)
%       };
%	\addplot+[no marks, domain=0.6:1.8, color={blue}, line width=1] (0.3,x); % an arc of R=0.3 from X = 0.6:1.8
   \end{smithchart}
   \end{tikzpicture} 
   \end{center}
   \end{textblock}
\end{frame}

%%%%%%%%%%%%%%%%%%%%%%%%%%%%%%%%%%%%%%%%%
%%%%%%%%%%%%%%%%%%%%%%%%%%%%%%%%%%%%%%%%%

\begin{frame}{Smith Charts: Measurements}
	\begin{textblock}{14.5}(0.5,3)
%	\only<1>{
%		\begin{tikzpicture}[scale=0.4]
%		    \begin{smithchart}[
%		        width=20cm,
%		        samples=200, line cap=round, smooth]
%		
%		    \end{smithchart}
%		
%		\end{tikzpicture}
%	} % end only 
	
	\begin{tikzpicture}[scale=0.4]
	    \begin{smithchart}[
	        width=20cm,
	        samples=200, line cap=round, smooth]
	
		\addplot[color=blue, thick, domain=0:30] (x,-1);
		\addplot[color=blue, thick, domain=0:30] (x,-0.9);
		\addplot[color=blue, thick, domain=0:30] (x,-0.8);
		\addplot[color=blue, thick, domain=0:30] (x,-0.7);
		\addplot[color=blue, thick, domain=0:30] (x,-0.6);
		\addplot[color=blue, thick, domain=0:30] (x,-0.5);
		\addplot[color=blue, thick, domain=0:30] (x,-0.4);
		\addplot[color=blue, thick, domain=0:30] (x,-0.3);
		\addplot[color=blue, thick, domain=0:30] (x,-0.2);
		\addplot[color=blue, thick, domain=0:30] (x,-0.1);
		\addplot[color=blue, thick, domain=0:30] (x,0.1);
		\addplot[color=blue, thick, domain=0:30] (x,-1);
		\addplot[color=blue, thick, domain=0:30] (x,0.9);
		\addplot[color=blue, thick, domain=0:30] (x,0.8);
		\addplot[color=blue, thick, domain=0:30] (x,0.7);
		\addplot[color=blue, thick, domain=0:30] (x,0.6);
		\addplot[color=blue, thick, domain=0:30] (x,0.5);
		\addplot[color=blue, thick, domain=0:30] (x,0.4);
		\addplot[color=blue, thick, domain=0:30] (x,0.3);
		\addplot[color=blue, thick, domain=0:30] (x,0.2);
		
	
		\addplot[color=red, thick, domain=30:-30, smooth, samples=400] (0.1,x); 	
		\addplot[color=red, thick, domain=30:-30, smooth] (0.2,x);
		\addplot[color=red, thick, domain=30:-30, smooth] (0.3,x);
		\addplot[color=red, thick, domain=30:-30, smooth] (0.4,x);
		\addplot[color=red, thick, domain=30:-30, smooth] (0.5,x);
		\addplot[color=red, thick, domain=30:-30, smooth] (0.6,x);
		\addplot[color=red, thick, domain=30:-30, smooth] (0.7,x);
		\addplot[color=red, thick, domain=30:-30, smooth] (0.8,x);
		\addplot[color=red, thick, domain=30:-30, smooth] (0.9,x);
		\addplot[color=red, thick, domain=30:-30, smooth] (1,x);
		\addplot[color=red, thick, domain=30:-30, smooth] (1.2,x);
		\addplot[color=red, thick, domain=30:-30, smooth] (1.4,x);
		\addplot[color=red, thick, domain=30:-30, smooth] (2,x);
		\addplot[color=red, thick, domain=30:-30, smooth] (3,x);
		\addplot[color=red, thick, domain=30:-30, smooth] (4,x);
		\addplot[color=red, thick, domain=30:-30, smooth] (5,x);
		\addplot[color=red, thick, domain=30:-30, smooth] (10,x);
		\addplot[color=red, thick, domain=30:-30, smooth] (20,x);  
		\draw[ultra thick, red,<-] (0.1,-0.25) -- (0.2,-0.4) node[above, rectangle, fill=white] {\Large{Constant Resistance}};
		 \draw[ultra thick, blue,<-] (1.8,0.9) -- (1.4,1.6) node[above, rectangle, fill=white] {\Large{Constant Reactance}};
	
		\path[draw=green!50!black, ultra thick] (1,0) circle (2cm);
		 \draw[ultra thick, green!50!black,<-] (0.6,0) -- (0.2,0.2) node[above, rectangle, fill=white] {\Large{Constant VSWR}};
	
		\node[rectangle, fill=white] at (0.4,-1) {\Huge{Capacitive}};		
		\node[rectangle, fill=white] at (0.4,1) {\Huge{Inductive}};
  	  \end{smithchart}

	\end{tikzpicture}


\end{textblock}
\begin{textblock}{5}(9.5,2.5)
	\begin{myblock}[0.95]{}
	
		\begin{itemize}
		
			\setlength\itemsep{12pt}
			\item Plot complex impedances
			\item Use circles of constant R \& X
			\item Measure $\Gamma_{\text{L}}$
			\item Magnitude and Phase
			\item 
		
		\end{itemize}
	
	\end{myblock}
\end{textblock}
\end{frame}
%\end{comment}
%%%%%%%%%%%%%%%%%%%%%%%%%%%%%%%%%%%%%%%%%%%
%%%%%%%%%%%%%%%%%%%%%%%%%%%%%%%%%%%%%%%%%%%

%\begin{frame}{Smith Charts: Official}
%
%\begin{textblock}{14.5}(0.5,2)
%
%\begin{center}
%	\includegraphics[scale=0.2]{ECE343/smith_chart}
%\end{center}
%
%\end{textblock}
%\end{frame}
%
%%%%%%%%%%%%%%%%%%%%%%%%%%%%%%%%%%%%%%%%%%%
%%%%%%%%%%%%%%%%%%%%%%%%%%%%%%%%%%%%%%%%%%%
%
%\begin{frame}{Smith Charts: Official}
%
%\begin{textblock}{14.5}(0.5,3)
%
%\begin{center}
%	\includegraphics[scale=0.75]{ECE343/color_smithchart}
%\end{center}
%
%\end{textblock}
%
%\end{frame}

{
%\setbeamercolor{background canvas}{bg=}
\includepdf[pages=1, landscape=false]{\mygraphicspath/ECE343/SmithChart_blank.pdf}
}

{
%\setbeamercolor{background canvas}{bg=}
\includepdf[pages=1, landscape=false]{\mygraphicspath/ECE343/colour_smith_chart.pdf}
}

%%%%%%%%%%%%%%%%%%%%%%%%%%%%%%%%%%%%%%%%%
%%%%%%%%%%%%%%%%%%%%%%%%%%%%%%%%%%%%%%%%%

\begin{frame}{Smith Charts: Components }

\begin{textblock}{10}(0.5,3)

		\begin{tikzpicture}[scale=0.4]
		    \begin{smithchart}[
		        width=20cm,
		        samples=50, line cap=round, smooth]
			%\addplot[ultra thick, red, mark=*, mark color=red] (2,-1.4) node[below, rectangle, fill=white]{\Large{$2-j1.4$}};
			\addplot[-<, red, ultra thick, domain=0.1:0.6] (0.5, x) node[left, rectangle, fill=white] {L};
			\addplot[-<, red, ultra thick, domain=-0.1:-0.6] (0.5, x) node[right, rectangle, fill=white] {C};
		    \end{smithchart}
		\end{tikzpicture}

\end{textblock}
\begin{textblock}{4.5}(10.5,3)

\begin{myblock}[0.95]{Series}
	\begin{itemize}
		\setlength\itemsep{14pt}
		\item Series \handoutgap{0.5in}{inductors} move \handoutgap{0.5in}{CW} along lines of constant resistance
	
		\item Series \handoutgap{0.5in}{capactiors} move \handoutgap{0.5in}{CCW} along lines of constant resistance
	\end{itemize}

\end{myblock}
\end{textblock}

\end{frame}

%%%%%%%%%%%%%%%%%%%%%%%%%%%%%%%%%%%%%%%%%%%
%%%%%%%%%%%%%%%%%%%%%%%%%%%%%%%%%%%%%%%%%%%

\begin{frame}{Smith Charts: Components }

\begin{textblock}{10}(0.5,3)

		\begin{tikzpicture}[scale=0.4]
		    \begin{smithchart}[
		        width=20cm,
		        samples=50, line cap=round, smooth]
			%\addplot[ultra thick, red, mark=*, mark color=red] (2,-1.4) node[below, rectangle, fill=white]{\Large{$2-j1.4$}};
			\addplot[-<, red, ultra thick, domain=0.1:0.6, mark=none] (0.5, x) node[left, rectangle, fill=white] {L};
			\addplot[-<, red, ultra thick, domain=-0.1:-0.6, mark=none] (0.5, x) node[right, rectangle, fill=white] {C};
		    \end{smithchart}
		\end{tikzpicture}

\end{textblock}
\begin{textblock}{4.5}(10.5,3)

\begin{myblock}[0.95]{Parallel}
	\begin{itemize}
		\setlength\itemsep{14pt}
		\item Parallel \handoutgap{0.5in}{inductors} move \handoutgap{0.5in}{CCW} along lines of constant conductance
	
		\item Parallel \handoutgap{0.5in}{capactiors} move \handoutgap{0.5in}{CW} along lines of constant conductance
	\end{itemize}

\end{myblock}
\end{textblock}

\end{frame}

%{
%%\setbeamercolor{background canvas}{bg=}
%\includepdf[pages=1, landscape=false]{ECE343/SmithChart_blank.pdf}
%}

%
%{
%%\setbeamercolor{background canvas}{bg=}
%\includepdf[pages=1, landscape=false]{ECE343/SmithChart_blank.pdf}
%}

%%%%%%%%%%%%%%%%%%%%%%%%%%%%%%%%%%%%%%%%%
%%%%%%%%%%%%%%%%%%%%%%%%%%%%%%%%%%%%%%%%%

\begin{frame}{Smith Charts: Plotting Impedances}

\begin{textblock}{10}(0.5,3)

		\begin{tikzpicture}[scale=0.4]
		    \begin{smithchart}[
		        width=20cm,
		        samples=200, line cap=round, smooth]
			%\addplot[ultra thick, red, mark=*, mark color=red] (2,-1.4) node[below, rectangle, fill=white]{\Large{$2-j1.4$}};
		    \end{smithchart}
		\end{tikzpicture}

\end{textblock}
\begin{textblock}{4.5}(10.5,3)

\begin{myblock}[0.95]{Plot these}

	\begin{itemize}
		\setlength\itemsep{8pt}
		\item $\myz{L} = 2-j1.4$ 
		
			$\longrightarrow$ (C or L?)
		
		\item $\myz{L} = 1.5+j0.5$ 
		
				$\longrightarrow$ (C or L?)
		
		\item $\myz{L} = 75+j60\ohms$  with \myz{0} = 50\ohms
		
		\item Short circuit
		
		\item Open circuit 
	\end{itemize}

\end{myblock}
\end{textblock}

\end{frame}

%%%%%%%%%%%%%%%%%%%%%%%%%%%%%%%%%%%%%%%%%%%%%%%%%%%%%%%%%%%%%%%%%%
%%%%%%%%%%%%%%%%%%%%%%%%%%%%%%%%%%%%%%%%%%%%%%%%%%%%%%%%%%%%%%%%%%
% LO 3.6
\begin{frame}{Impedance Matching: General Process}

	\begin{textblock}{14.5}(0.5,3)

	Match the combined load (\myz{L} and TX line) to \myz{s}.
		
		\vspace{1in}
		\scalebox{0.55}{
		\begin{circuitikz}
		\draw (0,0) to [sV, a^=\myv{s}] ++(0,3) 
				to[twoport, t=\myz{s}] ++(1,0)
				to[tline,o-o] ++(3.5,0)
				to[short] ++(0.5,0)
				to[twoport,t=\myz{L}]  ++(0,-3)
				to[short] ++(-0.5,0)
				to[tline,o-o] ++(-3.5,0)
				to[short] ++(-1,0);
	%	\draw[<->] (0.5,0.25) -- node[above] {$\ell$} ++(3,0); 
		\node[below] at (2,1.75) {\myz{0}};
		\node[below] at (0.75,0) {$-\ell$};
		\node[below] at (4.5,0) {$0$};
		\draw[->] (1.25,-0.5) --++(0,1) --++ (0.5,0);
		\node[below] at (1.25,-0.5) {\myz{in}};	
		\end{circuitikz}
		} % end scalebox
	\end{textblock}
	
	\begin{textblock}{11}(4.75,2.75)
	\begin{myblock}[0.95]{Process}
	
	\begin{small}
	\begin{enumerate}
	
		\item Normalize impedances to \myz{0}
		
		\item Find reflection coefficient ($\Gamma_L$)
		
		\item Holding the \handoutgap{0.75in}{magnitude} of $\Gamma_L$ constant, find the input impedance $z_{\text{in}} = \handoutgap{0.5in}{1 + jx}$ 
		
		$\implies |\Gamma_L | = \dfrac{| z_{\text{in}} - 1 |}{|z_{\text{in}} + 1|}$
		
		\item Determine capacitance or inductance from $X$. 
		
		$\implies Z_C = \dfrac{1}{j\omega C} \text{ or } Z_L = j\omega L$
		
		\item Find the length of TX line to make:
		
		$\implies z_{\text{in}} = 1+jX =  \dfrac{z_L + j\tanp{\beta \ell}}{1+jz_L\tanp{\beta\ell}}$
		

	\end{enumerate}
	\end{small}
	
	\end{myblock}
	\end{textblock}


	
\end{frame}


%%%%%%%%%%%%%%%%%%%%%%%%%%%%%%%%%%%%%%%%%%%%%%%%%%%%%%%%%%%%%%%%%%
%%%%%%%%%%%%%%%%%%%%%%%%%%%%%%%%%%%%%%%%%%%%%%%%%%%%%%%%%%%%%%%%%%

%\section{Review}
%\begin{frame}{Impedance Matching: General Process}
%
%	\begin{textblock}{14.5}(0.5,3)
%
%	Match the combined load (\myz{L} and TX line) to \myz{s}.
%		
%		\vspace{1in}
%		\scalebox{0.55}{
%		\begin{circuitikz}
%		\draw (0,0) to [sV, a^=\myv{s}] ++(0,3) 
%				to[twoport, t=\myz{s}] ++(1,0)
%				to[tline,o-o] ++(3.5,0)
%				to[short] ++(0.5,0)
%				to[twoport,t=\myz{L}]  ++(0,-3)
%				to[short] ++(-0.5,0)
%				to[tline,o-o] ++(-3.5,0)
%				to[short] ++(-1,0);
%	%	\draw[<->] (0.5,0.25) -- node[above] {$\ell$} ++(3,0); 
%		\node[below] at (2,1.75) {\myz{0}};
%		\node[below] at (0.75,0) {$-\ell$};
%		\node[below] at (4.5,0) {$0$};
%		\draw[->] (1.25,-0.5) --++(0,1) --++ (0.5,0);
%		\node[below] at (1.25,-0.5) {\myz{in}};	
%		\end{circuitikz}
%		} % end scalebox
%	\end{textblock}
%	
%	\begin{textblock}{11}(4.75,2.75)
%	\begin{myblock}[0.95]{Process}
%	
%	\begin{small}
%	\begin{enumerate}
%	
%		\item Normalize impedances to \myz{0}
%		
%		\item Find reflection coefficient ($\Gamma_L$)
%		
%		\item Holding the \handoutgap{0.75in}{magnitude} of $\Gamma_L$ constant, find the input impedance $z_{\text{in}} = \handoutgap{0.5in}{1 + jx}$ 
%		
%		$\implies |\Gamma_L | = \dfrac{z_{\text{in}} - 1}{z_{\text{in}} + 1}$
%		
%		\item Determine capacitance or inductance from $X$. 
%		
%		$\implies Z_C = \dfrac{1}{j\omega C} \text{ or } Z_L = j\omega L$
%		
%		\item Find the length of TX line to make:
%		
%		$\implies z_{\text{in}} = 1+jX =  \dfrac{z_L + j\tanp{\beta \ell}}{1+jz_L\tanp{\beta\ell}}$
%		
%
%	\end{enumerate}
%	\end{small}
%	
%	\end{myblock}
%	\end{textblock}
%
%
%	
%\end{frame}

%%%%%%%%%%%%%%%%%%%%%%%%%%%%%%%%%%%%%%%%%%%%%%%%%%%%%%%%%%%%%%%%%%
%%%%%%%%%%%%%%%%%%%%%%%%%%%%%%%%%%%%%%%%%%%%%%%%%%%%%%%%%%%%%%%%%%

\begin{frame}{Impedance Matching: Smith Chart Process}

	\begin{textblock}{14.5}(0.5,3)

	Match the combined load (\myz{L} and TX line) to \myz{s}.
		
		\vspace{1in}
		\scalebox{0.55}{
		\begin{circuitikz}
		\draw (0,0) to [sV, a^=\myv{s}] ++(0,3) 
				to[twoport, t=\myz{s}] ++(1,0)
				to[tline,o-o] ++(3.5,0)
				to[short] ++(0.5,0)
				to[twoport,t=\myz{L}]  ++(0,-3)
				to[short] ++(-0.5,0)
				to[tline,o-o] ++(-3.5,0)
				to[short] ++(-1,0);
	%	\draw[<->] (0.5,0.25) -- node[above] {$\ell$} ++(3,0); 
		\node[below] at (2,1.75) {\myz{0}};
		\node[below] at (0.75,0) {$-\ell$};
		\node[below] at (4.5,0) {$0$};
		\draw[->] (1.25,-0.5) --++(0,1) --++ (0.5,0);
		\node[below] at (1.25,-0.5) {\myz{in}};	
		\end{circuitikz}
		} % end scalebox
	\end{textblock}
	
	\begin{textblock}{11}(4.75,2.75)
	\begin{myblock}[0.95]{Process}
	
	\begin{small}
	\begin{enumerate}
	
		\item Normalize impedances to \myz{0}
		
		\item Draw \handoutgap{0.5in}{constant $\Gamma_L$ circle}
		
		\item Rotate towards \handoutgap{0.5in}{source} on circle until intersection with  \handoutgap{0.5in}{$z_s$ circle}
		
		\item Read \handoutgap{0.5in}{$z_{\text{in}}$}
		
		\item Determine capacitance or inductance from $X$. 
		
		$\implies Z_C = \dfrac{1}{j\omega C} \text{ or } Z_L = j\omega L$
		
		\item Measure TX line length ($\ell$) from \handoutgap{0.5in}{outer scale}
		

	\end{enumerate}
	\end{small}
	
	\end{myblock}
	\end{textblock}
	
\end{frame}


%%%%%%%%%%%%%%%%%%%%%%%%%%%%%%%%%%%%%%%%%%%%
%%%%%%%%%%%%%%%%%%%%%%%%%%%%%%%%%%%%%%%%%%%%
\section{Matching Example}
\begin{frame}{Impedance Matching Example}
\begin{textblock}{9}(0.5, 3)
	\includegraphics[scale=0.22]{ECE343/smithchart.eps}
\end{textblock}
\begin{textblock}{5.5}(10, 2.75)
\begin{myblock}[0.97]{Previous Example}
\begin{itemize}
	\setlength{\itemsep}{6pt}
	\item Match \myz{s}=50\ohms to \myz{L}=33.9+j17.9\ohms
	\item \myz{0}  = 50\ohms
	\item \handoutgap{1in}{Plot $z_L$}
	\item \handoutgap{1in}{Draw $\Gamma_L$ circle}
	\item \handoutgap{1in}{Rotate towards source}
	\item \handoutgap{1in}{Force Re$\lb z_{\text{in}}  \rb =z_s$}
	\item \handoutgap{1in}{Read $x$, unitize, find C/L}

\end{itemize}

\end{myblock}

\end{textblock}
\end{frame}
%\end{comment}
%%%%%%%%%%%%%%%%%%%%%%%%%%%%%%%%%%%%%%%%%%%
%%%%%%%%%%%%%%%%%%%%%%%%%%%%%%%%%%%%%%%%%%%
\begin{frame}{Impedance Matching Example}
%\begin{textblock}{14.5}(0.5,2.5)
%\begin{center}
%\begin{tikzpicture}[scale=0.35]
%
%\input{\mygraphicspath/ECE343/smithChartFoundation}
%
%   \begin{smithchart}[
%                      show origin,
%                      width=20cm,
%                     ]
%%		\addplot[mark=*, color=red] coordinates{(0.678,0.358)} node[above]{\myz{L}};
%%   \addplot[mark=none,line width=2]
%%       coordinates{
%%           (1, 0) (1, 0.1) (1,0.2) (1,0.3) (1,0.4) (1,0.5) (1,0.5)
%%       };
%%   \addplot[mark=none,line width=1, red]
%%       coordinates{
%%           %(0,-0.5)(1, 0) (0, 0.5) % pacman angle
%%           (1,0) (0.5,0.5)
%%       };
%%	\addplot+[no marks, domain=0.6:1.8, color={blue}, line width=1] (0.3,x); % an arc of R=0.3 from X = 0.6:1.8
%   \end{smithchart}
%   \end{tikzpicture} 
%   \end{center}
%   \end{textblock}
\begin{textblock}{9}(3.5, 3)
	\includegraphics[scale=0.22]{ECE343/smithchart.eps}
\end{textblock}
\end{frame}

%\begin{comment}
%%%%%%%%%%%%%%%%%%%%%%%%%%%%%%%%%%%%%%%%%%%%
%%%%%%%%%%%%%%%%%%%%%%%%%%%%%%%%%%%%%%%%%%%%
\section{Matching Practice}
\begin{frame}{Impedance Matching: Your Turn}
\begin{textblock}{9}(0.5, 3)
	%\scalebox{0.35}{
	%\begin{mysmithchart}
	%   \addplot[mark=none,line width=2]
	%       coordinates{
	%           (1, 0) (1, 0.1) (1,0.2) (1,0.3) (1,0.4) (1,0.5) (1,0.5)
	%       };
	%   \addplot[mark=none,line width=0.5]
	%       coordinates{
	%           (1, 0) (-0.3, 0)  % this one is not drawn outside!!!
	%       };
	%\end{mysmithchart}
	 %  } % end scalebox
	% \includesvg[scale=0.5]{ECE343/smithchart}
	\includegraphics[scale=0.22]{ECE343/smithchart.eps}
\end{textblock}
\begin{textblock}{5.5}(10, 2.75)
\begin{myblock}[0.97]{Example 2}
\begin{itemize}
	\setlength{\itemsep}{6pt}
	\item Match \myz{s}=50\ohms to \myz{L}=30-j120\ohms
	\item \myz{0}  = 75\ohms
	\item \handoutgap{1in}{Plot $z_L$}
	\item \handoutgap{1in}{Draw $\Gamma_L$ circle}
	\item \handoutgap{1in}{Rotate towards source}
	\item \handoutgap{1in}{Force Re$\lb z_{\text{in}}  \rb =1$}
	\item \handoutgap{1in}{Read $x$ and find C/L}

\end{itemize}

\end{myblock}

\end{textblock}
\end{frame}

%%%%%%%%%%%%%%%%%%%%%%%%%%%%%%%%%%%%%%%%%%%%
%%%%%%%%%%%%%%%%%%%%%%%%%%%%%%%%%%%%%%%%%%%%

\begin{frame}{Impedance Matching: Your Turn}
\begin{textblock}{9}(3.5, 3)
	\includegraphics[scale=0.22]{ECE343/smithchart.eps}
\end{textblock}
\end{frame}

%%%%%%%%%%%%%%%%%%%%%%%%%%%%%%%%%%%%%%%%%%%%
%%%%%%%%%%%%%%%%%%%%%%%%%%%%%%%%%%%%%%%%%%%%
\section{Parallel Components}
\begin{frame}{Impedance Matching: Parallel Component}
\begin{textblock}{9}(0.5, 3)
	\includegraphics[scale=0.22]{ECE343/smithchart.eps}
\end{textblock}
\begin{textblock}{5.5}(10, 2.75)
\begin{myblock}[0.97]{Previous Example}
\begin{itemize}
	\setlength{\itemsep}{6pt}
	\item Use a parallel component
	\item \handoutgap{1in}{Plot $z_L$}
	\item \handoutgap{1in}{$z_L \rightarrow y_L $}
	\item \handoutgap{1in}{Draw $\Gamma_L$ circle}
	\item \handoutgap{1in}{Rotate towards source}
	\item \handoutgap{1in}{Force Re$\lb y_{\text{in}}  \rb =y_s$}
	\item \handoutgap{1in}{Read $b$, unitize, and find C/L}

\end{itemize}

\end{myblock}

\end{textblock}
\end{frame}

%%%%%%%%%%%%%%%%%%%%%%%%%%%%%%%%%%%%%%%%%%%%
%%%%%%%%%%%%%%%%%%%%%%%%%%%%%%%%%%%%%%%%%%%%

\begin{frame}{Impedance Matching: Parallel Component}
\begin{textblock}{9}(3.5, 3)
	\includegraphics[scale=0.22]{ECE343/smithchart.eps}
\end{textblock}
\end{frame}


%%%%%%%%%%%%%%%%%%%%%%%%%%%%%%%%%%%%%%%%%%%%
%%%%%%%%%%%%%%%%%%%%%%%%%%%%%%%%%%%%%%%%%%%%
\section{Stubs}
\begin{frame}{Impedance Matching: TX Line Stubs}
\begin{textblock}{14.5}(0.5, 3.5)
\scalebox{0.65}{
\begin{circuitikz}
	\draw (0,0) to[short,o-o] ++ (3,0)
				to[short,-o] ++(4,0)
				to[short] ++(0.5,0)
				to[twoport, t=\myz{L}, a^={\myz{L}\equals 33.9+j17.9\ohms}] ++(0,-2)
				to[short] ++(-0.5,0)
				to[short,o-] ++(-4,0)
				to[short,o-o] ++(-3,0);
	\draw[<->] (3,0.5) -- node[above] {$d$} ++(4,0); 
	\draw(1.5,-1) node{\myz{s} = 50\ohms};
	\draw(5,-1) node{\myz{0,2} = 50\ohms};
	\draw[<->] (2.8,-2.2) -- node[label={[label distance=8pt]195:$\ell$}] {} ++(1.5,-2);
	\draw(3,0) to[short] ++(1.75,-2.25) to[open,o-o] ++(0,-2) to[short] ++(-1.75,2.25);
	\draw(5.5,-3) node{\myz{0,stub} = 100\ohms};
\end{circuitikz}
} % end scalebox

\scalebox{0.65}{
\begin{circuitikz}
	\draw (0,0) to[short,o-o] ++ (3,0)
				to[short,-o] ++(4,0)
				to[short] ++(0.5,0)
				to[twoport, t=\myz{L}, a^={\myz{L}\equals 33.9+j17.9\ohms}] ++(0,-2)
				to[short] ++(-0.5,0)
				to[short,o-] ++(-4,0)
				to[short,o-o] ++(-3,0);
	\draw[<->] (3,0.5) -- node[above] {$d$} ++(4,0); 
	\draw(1.5,-1) node{\myz{s} = 50\ohms};
	\draw(5,-1) node{\myz{0,2} = 50\ohms};
	\draw[<->] (2.8,-2.2) -- node[label={[label distance=8pt]195:$\ell$}] {} ++(1.5,-2);
	\draw(3,0) to[short] ++(1.75,-2.25) to[short] ++(0,-2) to[short] ++(-1.75,2.25);
	\draw(5.5,-3) node[right]{\myz{0,stub} = 50\ohms};
\end{circuitikz}
} % end scalebox

\end{textblock}
\end{frame}


%%%%%%%%%%%%%%%%%%%%%%%%%%%%%%%%%%%%%%%%%%%%
%%%%%%%%%%%%%%%%%%%%%%%%%%%%%%%%%%%%%%%%%%%%

\begin{frame}{Impedance Matching: Shorted Stub}
\begin{textblock}{9}(0.5, 3)
	\includegraphics[scale=0.22]{ECE343/smithchart.eps}
\end{textblock}
\begin{textblock}{5.5}(10, 2.75)
\begin{myblock}[0.97]{Previous Example}
\begin{itemize}
	\setlength{\itemsep}{6pt}
	\item Match \myz{s}=50\ohms to \myz{L}=33.9+j17.9\ohms
	\item \handoutgap{1in}{$z_L \rightarrow y_L$}
	\item \handoutgap{1in}{Draw $\Gamma_L$ circle}
	\item \handoutgap{1in}{Rotate towards source}
	\item \handoutgap{1in}{Force Re$\lb z_{\text{in}}  \rb =1$}
	\item \handoutgap{1in}{Read $x$, find on outer circle}
	\item \handoutgap{1in}{Find $\ell$ from $y_{\text{sc}}$ or $y_{\text{sc}}$ to $x$}

\end{itemize}

\end{myblock}

\end{textblock}
\end{frame}


%%%%%%%%%%%%%%%%%%%%%%%%%%%%%%%%%%%%%%%%%%%%
%%%%%%%%%%%%%%%%%%%%%%%%%%%%%%%%%%%%%%%%%%%%

\begin{frame}{Impedance Matching: Shorted Stub}
\begin{textblock}{9}(3.5, 3)
	\includegraphics[scale=0.22]{ECE343/smithchart.eps}
\end{textblock}
\end{frame}


%%%%%%%%%%%%%%%%%%%%%%%%%%%%%%%%%%%%%%%%%%%%
%%%%%%%%%%%%%%%%%%%%%%%%%%%%%%%%%%%%%%%%%%%%

\begin{frame}{Impedance Matching: Open Stub - Your Turn}
\begin{textblock}{9}(0.5, 3)
	%\scalebox{0.35}{
	%\begin{mysmithchart}
	%   \addplot[mark=none,line width=2]
	%       coordinates{
	%           (1, 0) (1, 0.1) (1,0.2) (1,0.3) (1,0.4) (1,0.5) (1,0.5)
	%       };
	%   \addplot[mark=none,line width=0.5]
	%       coordinates{
	%           (1, 0) (-0.3, 0)  % this one is not drawn outside!!!
	%       };
	%\end{mysmithchart}
	 %  } % end scalebox
	% \includesvg[scale=0.5]{ECE343/smithchart}
	\includegraphics[scale=0.22]{ECE343/smithchart.eps}
\end{textblock}
\begin{textblock}{5.5}(10, 2.75)
\begin{myblock}[0.97]{Example 2}
\begin{itemize}
	\setlength{\itemsep}{6pt}
	\item Match \myz{s}=50\ohms to \myz{L}=33.9+j17.9\ohms
	\item \handoutgap{1in}{$z_L \rightarrow y_L$}
	\item \handoutgap{1in}{Draw $\Gamma_L$ circle}
	\item \handoutgap{1in}{Rotate towards source}
	\item \handoutgap{1in}{Force Re$\lb z_{\text{in}}  \rb =1$}
	\item \handoutgap{1in}{Read $x$, find on outer circle}
	\item \handoutgap{1in}{Find $\ell$ from $y_{\text{sc}}$ or $y_{\text{sc}}$ to $x$}

\end{itemize}

\end{myblock}

\end{textblock}
\end{frame}


%%%%%%%%%%%%%%%%%%%%%%%%%%%%%%%%%%%%%%%%%%%%
%%%%%%%%%%%%%%%%%%%%%%%%%%%%%%%%%%%%%%%%%%%%

\begin{frame}{Impedance Matching: Open Stub - Your Turn}
\begin{textblock}{9}(3.5, 3)
	\includegraphics[scale=0.22]{ECE343/smithchart.eps}
\end{textblock}
\end{frame}

%\end{comment}

